% Section 1. Input from paper.tex; every number traces to sources.md and is asserted by check_numbers.py.
\providecommand{\lm}{\texttt{local\_mean\_28}}
\providecommand{\qs}{QScore9}

\section{Introduction}

A daily count of new stars on a GitHub repository, of downloads of an npm package or of 311 service
requests for one problem in one community board is zero on most days for most series, and small when it is
not. A forecast of such a series has two parts: how likely a day is to be empty, and how large the count
will be when it is not. The zero day is the event of interest for these counts, not a nuisance: no new star,
no download and no complaint is what a maintainer, a registry operator or a city agency watches for, and a
forecaster that cannot say how likely it is cannot be used to alert on a change in activity. Zero-shot time-series foundation models forecast such a series without fitting, as nine quantiles. The practical question is whether one of them is worth running on
a sparse count series at all, or whether a trailing mean does as well, and where any value it adds sits.

This paper's answer is that the value is in the size, not the occurrence. On sparse daily counts the frozen
models' gain over a 28-day local mean comes from the days that are not zero, where the local mean
overpredicts. Their own probability of a zero day is at least as good as a context-only estimate wherever
their nine quantiles can express one: on dense series in all three settings, where substituting a context
estimate for the model's zero mass makes the forecast worse, and on sparse NYC 311 series, where the two are
about level. The one place where a context estimate beats the model's own zero mass on the typical series is
the sparse npm stratum at a week, a secondary reading on one domain and one horizon. And on series that are
zero nine days in ten the models' nine deciles are all zero: the rule that completes them into a count law
then sets the probability of a positive day between 0.0607 and 0.1 whatever the model forecast, the quantile
score ties, and nothing about occurrence can be read from the models there.

The three studies are one design run in three steps, each registered before its target data was read and
each prompted by a result of the one before. Study 1 compared the raw forecasts of frozen Chronos-2, TimesFM
3 and TimesFM 2.5 with a 28-day local mean and weekly AutoETS on 300 GitHub repositories sampled from a frame
of 148,550, and found the models ahead on the mean and behind on the typical repository, with a post-hoc
accounting that put the mean gain on positive days and a small loss on zero days. Study 2 asked whether that
loss could be taken out by replacing each model's zero mass with a 28-day context-only occurrence
probability while keeping its completed positive law, on 209 hash-sampled npm download series and a
540-series synthetic grid; its pooled families answered no, and its declared secondary strata showed the
typical sparse npm series favoring the substitution at a week. Study 3 registered that regime reading as a
per-stratum family of 36 tests on 300 NYC 311 service-request series with a dynamic occurrence estimate.

The paper makes three contributions: the record, a registered, population-sampled evaluation of frozen
forecasting models on three public count sources, with the protocol's registered bytes verifiable by hash,
every score recomputed from the saved forecasts with no disagreement, the mean and the typical series side by
side everywhere, and every number traced to a saved file; the occurrence result by regime, with the model's
positive law beating a shifted-Poisson count law wherever that law is misspecified and losing to it where it
is the generator's own; and the interface limit, that on very sparse series the registered construction,
nine deciles completed into a count law and scored by a quantile score, cannot ask the occurrence question,
because the completion fixes the models' occurrence probability and the score cannot see a difference
smaller than one decile.  The paper does not claim that the models forecast bursts, that any result transfers to launches or to the days after a release, that a download is a human adoption or a request a measure of the condition it reports, or that it is the first to evaluate frozen foundation models on sparse counts, to wrap one in a hurdle, or to show aggregate and per-series rankings disagreeing (Section~\ref{sec:related}). Section~\ref{sec:design} gives the
design once; Sections~\ref{sec:raw} to~\ref{sec:positive} answer the three questions in turn;
Section~\ref{sec:discussion} reads the studies together; Sections~\ref{sec:limitations} and~\ref{sec:repro}
give the limitations and what is released.
